\subsection{Segmentation Tree}

\subsubsection{Seg Normal}

\inccode{estrut_dados/seg.cpp}

\subsubsection{Memória eficiente}

Para fazer a seg usar $2n$ de memória em vez de $4n$, basta mudar como se calcula os índices:

\inccode{estrut_dados/seg-mem-efic.cpp}

Não notei nenhuma diferença na performance nos testes feitos.

\subsubsection{Lazy}

\inccode{estrut_dados/seg-lazy.cpp}

\subsubsection{Persistent}

Vértices apontam uns aos outros. Cada update aloca um vértice para cada um modificado, e cria uma nova raiz.

\inccode{estrut_dados/seg-persistent.cpp}

\subsubsection{Esparsa}

Cria os vértices sob demanda. Cuidado com uso de memória (talvez tenha que por os nós em um vetor, e salvar o índice em vez de ponteiros. Isso economiza 12 bytes por nó).

\inccode{estrut_dados/seg-esparsa.cpp}

